% Page budget: 1.55 pages | Owns: negation, OBC derivation, additional-state scope, implementation, and true-parameter condition.
% WRITING GUIDE
% Purpose: Explain negation, derive the final state-level orthogonal-by-construction map, and delimit its interpretability guarantee.
% Start here: Decisions D-185, D-190, D-192, D-194, D-198, D-202, D-203, D-222, D-229, and D-232 in docs/decisions.md.
% Supporting material: Current state-level OBC derivations under scripts/gantry/orthogonal-by-construction/documentation and their thesis-readiness guidance.
% Mandatory papers: Read Györök's projection-regularization paper and orthogonal-by-construction paper in full before writing; keep their guarantees and assumptions separate.
% Research-plan anchor: Preserve Aspect 3's goal of protecting physical interpretability; state clearly how the final construction differs from the proposed regularisation.
% Current decisions: Reduced identifiable coordinates, data-derived reference tuples, a tangent-only one-step physical-state projection, per-epoch refresh, and no affine comparison arm.
% Choices to justify: Protected parameter coordinates, reference-set construction, projected rows, rank tolerance, exclusion of the affine offset, refresh frequency, and treatment of additional states.
% Implementation audit: Read model_augmentation/fit_systems/obc.py, scripts/gantry/gantry_dynamic/obc_gantry.py, obc_diagnostics.py, and the OBC attachment in model.py before fixing the derivation.
% Reference check: Verify the original projection method and separate its claims from the state-level, nonlinear-parameter, LPV, and additional-state extensions developed here.
% Cautions: docs/orthogonal-projection-plan.md describes an older soft-penalty route; one-step state orthogonality is not trajectory-output orthogonality or guaranteed true recovery.
% Planned figures: New geometric projection figure with a physical and additional-state partition inset; show the remaining uniqueness limitation as clearly as the protected direction.
% Section is finished when: The protected space, coordinate frame, gradient treatment, added-state scope, and true-recovery condition are explicit.
\section{Interpretability by Orthogonal Construction}\label{sec:obc}

The augmented model of Section~\ref{sec:aug_structure} leaves the split between
$f_{\mathrm{base}}$ and $f_{\mathrm{aug}}$ non-unique.
We make it unique with the orthogonal-by-construction (OBC) augmentation of
Gy\"or\"ok et al.~\cite{gyorok2026obc}.
That method is developed for input-output models that are linear in their
parameters, and we adapt it to the self-scheduled state equation of the gantry.
Section~\ref{sec:obc_negation} identifies the direction that makes the split
non-unique, Section~\ref{sec:obc_map} constructs the orthogonal correction, and
Section~\ref{sec:obc_scope} states what the construction guarantees and when it
recovers the true parameters.

\subsection{Negation of the Physical Parameters}\label{sec:obc_negation}
Joint estimation fits $\theta_{\mathrm{base}}$ and $\theta_{\mathrm{aug}}$
together, and $f_{\mathrm{aug}}$ writes into the same physical rows as
$f_{\mathrm{base}}$ in \eqref{eq:aug_transition}.
Different parameter pairs can therefore realise the same one-step
map~\cite[Ex.~2]{gyorok2026obc}, \cite[Ex.~1]{gyorok2025l4dc}.
The directions in which $\theta_{\mathrm{base}}$ changes this map are the
columns of the transition sensitivity
\begin{equation}
 \Phi_{\bar\theta}(\tilde x,u)
 =\left.\frac{\partial f_{\mathrm{base}}(\theta_{\mathrm{base}},\tilde x,u)}
   {\partial\theta_{\mathrm{base}}}\right|_{\theta_{\mathrm{base}}=\bar\theta_{\mathrm{base}}}
   \in\R^{6\times10},
\label{eq:obc_sens}
\end{equation}
where $\bar\theta_{\mathrm{base}}$ is the expansion point.

A learned write along these directions cancels a parameter displacement
$\delta\theta$ to first order:
\begin{equation}
\begin{aligned}
 &f_{\mathrm{base}}(\bar\theta_{\mathrm{base}}+\delta\theta,\tilde x,u)
 +\bigl[f_{\mathrm{aug}}(\theta_{\mathrm{aug}},\hat x,u)
 -\Phi_{\bar\theta}(\tilde x,u)\,\delta\theta\bigr]\\
 &\quad\approx f_{\mathrm{base}}(\bar\theta_{\mathrm{base}},\tilde x,u)
 +f_{\mathrm{aug}}(\theta_{\mathrm{aug}},\hat x,u).
\end{aligned}
\label{eq:obc_negation}
\end{equation}
The data cannot attribute this component to either part.
We call this negation.
Under negation, the physical estimate depends on the
initialisation~\cite[Sec.~5.1]{gyorok2026obc}.

For the gantry, differentiating \eqref{eq:eom} with respect to the $j$-th
combination at fixed state and input gives
\begin{equation}
 \frac{\partial\ddot q}{\partial\theta_{\mathrm{base},j}}
 =-M(Y)^{-1}\Bigl[\frac{\partial M(Y)}{\partial\theta_{\mathrm{base},j}}\,\ddot q
 +\frac{\partial C}{\partial\theta_{\mathrm{base},j}}\,\dot q
 +\frac{\partial K}{\partial\theta_{\mathrm{base},j}}\,q\Bigr].
\label{eq:obc_gantry_sens}
\end{equation}
The inertia combinations act through $\ddot q$, the damping combinations
through $\dot q$, and the stiffness combination through $q$.
A negating write therefore acts as an added inertia, damping or stiffness
force.
Through $M(Y)^{-1}$, these directions depend on the payload position, so the
space to be protected is scheduled by $Y$.
The RK4 step of Section~\ref{sec:aug_structure} carries
\eqref{eq:obc_gantry_sens} into \eqref{eq:obc_sens}.

The sensitivity is taken with respect to the ten combinations
$\theta_{\mathrm{base}}$ of Section~\ref{sec:params}.
With respect to the fourteen raw parameters it would have four null
directions, along which no data decide the estimate, with or without
augmentation.
The combinations remove this structural ambiguity, so the remaining
non-uniqueness is negation alone.
Training updates $\theta_{\mathrm{base}}$ in logarithmic and bounded
coordinates (Section~\ref{sec:params}).
This coordinate map has an invertible Jacobian, so the column space of
\eqref{eq:obc_sens} is the same in both coordinates.

A parameter regulariser does not remove this direction.
The cost term of~\cite[Sec.~5.2, Eq.~(26)]{hoekstra2026lfr} penalises the
distance of $\theta_{\mathrm{base}}$ to a priori values.
Along \eqref{eq:obc_negation} the data fit is flat to first order, so this
term alone selects the estimate.
In this work the a priori values would be the detuned start
(Section~\ref{sec:setup}).

\subsection{Orthogonal-by-Construction State Update}\label{sec:obc_map}
Gy\"or\"ok et al.\ earlier penalised the component of the learned output that
lies in the range of the baseline~\cite[Eq.~(13)]{gyorok2025l4dc}.
The penalty weight trades data fit against orthogonality and must be
tuned~\cite[Sec.~1]{gyorok2026obc}.
The OBC parametrisation removes this weight.
For an input-output baseline $\phi(x_k)\theta_b$ that is linear in its
parameters, it fits the stacked network output to the stacked regressor $\Phi$
and subtracts the fit~\cite[Eqs.~(8), (10), (11), (13)]{gyorok2026obc}:
\begin{equation}
\begin{aligned}
 \theta_{\mathrm{aux}}&=(\Phi^\top\Phi)^{-1}\Phi^\top F^{\mathrm{ANN}}_{\theta_a},\\
 \widetilde F^{\mathrm{ANN}}_{\theta_a}&=F^{\mathrm{ANN}}_{\theta_a}-\Phi\,\theta_{\mathrm{aux}},\\
 0&=\Phi^\top\widetilde F^{\mathrm{ANN}}_{\theta_a},\\
 \hat y_k&=\phi(x_k)\theta_b+f^{\mathrm{ANN}}_{\theta_a}(x_k)-\phi(x_k)\,\theta_{\mathrm{aux}},
\end{aligned}
\label{eq:obc_gyorok}
\end{equation}
where $x_k$ collects lagged inputs and outputs, $F^{\mathrm{ANN}}_{\theta_a}$
stacks the network outputs on the training data, and the last line is the
predictor with $\theta_{\mathrm{aux}}$ fixed after training.
The third line holds for every $\theta_a$~\cite[Lemma~3]{gyorok2026obc}.
We therefore use this construction rather than the penalty.

We apply \eqref{eq:obc_gyorok} to the physical rows of
\eqref{eq:aug_transition}, with the sensitivity \eqref{eq:obc_sens} as
regressor.
$f_{\mathrm{base}}$ is rational in $\theta_{\mathrm{base}}$ through $M(Y)^{-1}$
and has no linear-in-parameters form.
The regressor is therefore its tangent at $\bar\theta_{\mathrm{base}}$, as
in~\cite[Eq.~(15)]{gyorok2025l4dc}.
The tangent defines only the protected space.
The rollout keeps the nonlinear $f_{\mathrm{base}}$.
Unlike~\cite[Eqs.~(17), (18)]{gyorok2025l4dc}, the offset column of the
expansion is not protected.
By \eqref{eq:obc_negation}, a negating write lies in $\ran\Phi_{\bar\theta}$.
The construction removes the part of any write in this range, including the
part of a write along the offset.
The remainder cannot be produced by a change of $\theta_{\mathrm{base}}$ to
first order.

The regressor is evaluated at states, but only positions are measured.
Gy\"or\"ok et al.\ then evaluate it at states simulated with the
baseline~\cite[Sec.~3]{gyorok2025l4dc}.
An open-loop baseline simulation drifts on the free axes $X$ and $Y$
(Section~\ref{sec:aug_closed_loop}).
We therefore build fixed reference tuples from the measured positions of every
valid training sample,
\begin{equation}
 q_i=P^{-\top}y_i,\qquad
 \tilde x_i=\col(q_i,\dot q_i),\qquad
 \bar x_i=0,\qquad i=1,\dots,N,
\label{eq:obc_tuples}
\end{equation}
where $\dot q_i$ is the fourth-order central difference of $q$ within each
record and $u_i$ is the recorded actuator force.
Each tuple evaluates $M(Y_i)$ at its measured payload position, so the stack
covers the scheduling range of the training data.
The recorded input already contains the feedback action, so the controller
does not enter the tuples.
The tuples are normalised as the training signals
(Section~\ref{sec:aug_structure}).
A fixed auxiliary evaluation set is permitted
by~\cite[Sec.~3.2]{gyorok2026obc}.

On these tuples, the construction \eqref{eq:obc_gyorok} becomes
\begin{equation}
\begin{aligned}
 \Phi_{\bar\theta}(\tilde X,U)
 &=\col\bigl(\Phi_{\bar\theta}(\tilde x_1,u_1),\dots,
   \Phi_{\bar\theta}(\tilde x_N,u_N)\bigr)\in\R^{6N\times10},\\
 F_{\mathrm{aug}}
 &=\col\bigl(f_{\mathrm{aug}}(\theta_{\mathrm{aug}},\hat x_1,u_1),\dots,
   f_{\mathrm{aug}}(\theta_{\mathrm{aug}},\hat x_N,u_N)\bigr),\\
 \theta_{\mathrm{aux}}&=\Phi_{\bar\theta}(\tilde X,U)^{+}F_{\mathrm{aug}},\\
 0&=\Phi_{\bar\theta}(\tilde X,U)^\top
   \bigl(F_{\mathrm{aug}}-\Phi_{\bar\theta}(\tilde X,U)\,\theta_{\mathrm{aux}}\bigr),\\
 0&=\Phi_{\bar\theta}(\tilde X,U)^\top\,
   \frac{\partial\bigl(F_{\mathrm{aug}}-\Phi_{\bar\theta}(\tilde X,U)\,
   \theta_{\mathrm{aux}}\bigr)}{\partial\theta_{\mathrm{aug}}},
\end{aligned}
\label{eq:obc_ours}
\end{equation}
with $\hat x_i=\col(\tilde x_i,\bar x_i)$.
The stack must have full column rank, the counterpart
of~\cite[Assumption~1]{gyorok2026obc}.
Its rank is checked at every rebuild against the standard numerical tolerance
of its singular values.
Section~\ref{sec:results} reports the rank for the thesis data.
The fourth line is~\cite[Lemma~3]{gyorok2026obc}.
It holds for every $\theta_{\mathrm{aug}}$, and the stack does not depend on
$\theta_{\mathrm{aug}}$, so differentiating it gives the last line.
The network gradient of the corrected stack thus has no component in the
protected range.
This requires the gradient to pass through $\theta_{\mathrm{aux}}$.
With a detached coefficient, the gradient is not projected.

Gy\"or\"ok et al.\ either update the expansion point at every objective
evaluation or fix it at the nominal
parameters~\cite[Sec.~4]{gyorok2025l4dc}.
Updating rebuilds the full stack at every evaluation.
A fixed nominal point does not apply here, because training starts from
detuned parameters (Section~\ref{sec:setup}).
We rebuild the stack at the current $\theta_{\mathrm{base}}$ at the start of
every epoch and hold it fixed within the epoch.
No gradient therefore reaches $\theta_{\mathrm{base}}$ through the stack.
The coefficient $\theta_{\mathrm{aux}}$ is recomputed once per objective
evaluation, from the current network on the fixed tuples.
Fitting it on the rollout instead would make the correction at step $k$ depend
on later rollout states.

The corrected augmented model is
\begin{equation}
\begin{aligned}
 \tilde x_{k+1}
 &=f_{\mathrm{base}}(\theta_{\mathrm{base}},\tilde x_k,\hat u_k)
   +f_{\mathrm{aug}}(\theta_{\mathrm{aug}},\hat x_k,\hat u_k)
   -\Phi_{\bar\theta}(\tilde x_k,\hat u_k)\,\theta_{\mathrm{aux}},\\
 \bar x_{k+1}
 &=g_{\mathrm{aug}}(\theta_{\mathrm{aug}},\hat x_k,\hat u_k),\\
 \hat y_k&=C_n\tilde x_k,
\end{aligned}
\label{eq:obc_model}
\end{equation}
where $\hat u_k$ is the closed-loop input of \eqref{eq:aug_controller}.
At each rollout step, the correction evaluates the sensitivity at the epoch's
expansion point and applies the current objective's coefficient, as
in~\cite[Eq.~(13)]{gyorok2026obc}.
The baseline has zero rows for the additional states, so $g_{\mathrm{aug}}$ is
not corrected.
For prediction after training, the stack and $\theta_{\mathrm{aux}}$ are
rebuilt once at the selected parameters and then fixed.

\subsection{Scope and Condition for True Parameter Recovery}\label{sec:obc_scope}
The construction guarantees the identity in the fourth line of
\eqref{eq:obc_ours} and nothing more.
The identity is local in $\theta_{\mathrm{base}}$, because it protects the
tangent at $\bar\theta_{\mathrm{base}}$.
It holds for the stacked one-step map over the tuples, not at each tuple.
It uses the Euclidean inner product in normalised state coordinates, so the
state normalisation sets its metric.
The tuples lie on $\bar x=0$, so the coefficient does not see network outputs
that act only through nonzero additional states.
The additional-state rows are never corrected, and $g_{\mathrm{aug}}$ reaches
the physical rows one step later through $f_{\mathrm{aug}}$.
The construction therefore does not make the output trajectories of the two
components orthogonal, which is the level at which \eqref{eq:aug_loss} fits
the model.
The consistency and zero-covariance results
of~\cite[Thms.~16, 18]{gyorok2026obc} assume a predictor that is linear in the
baseline parameters and do not transfer.

A unique split is not yet the true split.
Let $\theta^*_{\mathrm{base}}$ be the true parameters and $\Delta^*$ the true
one-step discrepancy on the tuples,
\begin{equation}
 \Delta^*=\col\bigl(\tilde x_{i+1}-f_{\mathrm{base}}(\theta^*_{\mathrm{base}},
 \tilde x_i,u_i)\bigr)_{i=1}^{N},
\label{eq:obc_delta}
\end{equation}
where $\tilde x_{i+1}$ is the recorded successor, built with the same stencil
and normalisation as $\tilde x_i$.
Three assumptions are made.
The data are noise-free.
$f_{\mathrm{base}}$ is affine in $\theta_{\mathrm{base}}$ between
$\theta^*_{\mathrm{base}}$ and the estimate $\hat\theta_{\mathrm{base}}$.
The corrected one-step map reproduces every recorded successor, the
counterpart of~\cite[Assumption~6]{gyorok2026obc}.
Multiplying this one-step identity by $\Phi^\top$ removes the corrected write
by \eqref{eq:obc_ours} and gives
\begin{equation}
 \hat\theta_{\mathrm{base}}-\theta^*_{\mathrm{base}}=\Phi^{+}\Delta^*,
 \qquad
 \hat\theta_{\mathrm{base}}=\theta^*_{\mathrm{base}}
 \;\Longleftrightarrow\;
 \Phi^\top\Delta^*=0,
\label{eq:obc_bias}
\end{equation}
where $\Phi$ is the stack of \eqref{eq:obc_ours}, which the affine assumption
makes independent of the expansion point.
This is our state-level counterpart of~\cite[Cond.~4,
Eq.~(21)]{gyorok2026obc}.
Training minimises the closed-loop multistep error \eqref{eq:aug_loss}, so
\eqref{eq:obc_bias} predicts the parameter error of a one-step fit, not that
of the trained model.

The condition $\Phi^\top\Delta^*=0$ concerns the sum over all tuples, so it can
hold through cancellation where no single term vanishes.
Unless the missing dynamics are orthogonal to the regressor by structure, it
requires designed excitation~\cite[Sec.~4.1, Remark~5]{gyorok2026obc}.
Whether the thesis data satisfy it is measured by the in-range fraction of the
discrepancy,
\begin{equation}
 \rho^*=\frac{\norm{\Phi\Phi^{+}\Delta^*}}{\norm{\Delta^*}},
\label{eq:obc_overlap}
\end{equation}
with $\Phi$ evaluated at $\theta^*_{\mathrm{base}}$.
It is zero exactly when the recovery condition holds.
It requires $\theta^*_{\mathrm{base}}$ and is therefore available in
simulation only.
Section~\ref{sec:results} tests whether the split is start-independent and
relates the parameter error to $\rho^*$.
